\newpage
\section{Geometrie in $\mathbf{\mathbb{R}^n}$}
\label{2.4}

\begin{definition}{Skalarprodukt}
    Sei $V$ ein Vektorraum über $K\in\{\mathbb{R},\mathbb{C}\}$. Eine Abbildung $(\cdot,\cdot):V\times V\to K$ heißt Skalarprodukt, wenn gilt:
    \[
        (x,x)\ge 0,\ (x,x)=0\Rightarrow x=0,\qquad
        (x,y)=(y,x),\qquad
        (ax_1+bx_2,y)=a(x_1,y)+b(x_2,y).
    \]
\end{definition}

\begin{korollar}
    Sei $(X,d)$ ein metrischer Raum und $K\subset X$ kompakt. Ist $A\subset K$ abgeschlossen, so ist $A$ kompakt.
\end{korollar}


\begin{definition}[deflabel=ortgoOrtho]{Orthogonalität und Orthonormalsystem}
    Die Vektoren $x,y\in \K^n$ heißen \emph{orthogonal}, falls $(x,y)_2=0$. Eine Menge von Vektoren 
    $$
        \{a(1),\dots,a(m)\}, a^{(i)} \not= 0,a^{(i)} \in \K, i=1,\dots,m
    $$ 
    heißt \emph{Orthogonalsystem}, wenn 
    $$
        \underbrace{(a^{(k)},a^{(l)})=0}_{\text{paarweise orthogonal}} \quad \text{für } k\not= l \, .
    $$
    Ist zusätzlich $(a(k),a(k))_2=1$, so heißt es \emph{Orthonormalsystem}.

    Ist $m=n$, so handelt es sich um eine \emph{Orthogonalbasis}, bzw. wenn $(a(k),a(k))_2=1$ um eine \emph{Orthononalbasis}.
\end{definition}

\begin{bemerkungbox}
    Die orthogonalen Vektoren nach \cdef{ortgoOrtho} sind linear unabhänig:
    $$
    % \begin{aligned}
    %     \sum_{k=1}^{m} c_k\,a^{(k)} = 0 &\xLeftrightarrow[\text{mit a^{(l)}}]{\text{Skal.Prod.}} \sum_{k=1}^{m} c_k\left(a^{(k), a^{(l)}}\right) 
    %     \xlongequal{a^{(i)} \text{ paarweise orthog.}} = c_l \underbrace{(a^{(l)}, a^{(l)})}_{\not=0, \text{ für } a^{(l)} \not=0} = 0 \\
    %         &\Longleftrightarrow c_l = 0, \quad l = 1,\cdots,m
    % \end{aligned}
    \begin{aligned}
        \sum_{k=1}^{m} c_k\,a^{(k)} = 0
        &\xLeftrightarrow[\text{mit } a^{(l)}]{\text{Skal.Prod.}}
        \sum_{k=1}^{m} c_k\,(a^{(k)}, a^{(l)}) \\
        &\xlongequal[\text{orthogonal}]{a^{(i)} \text{ paarweise}}
        c_l \underbrace{(a^{(l)}, a^{(l)})}_{\neq 0,\; a^{(l)}\neq 0}
        = 0 \\
        &\Longleftrightarrow c_l = 0,\quad l=1,\dots,m.
    \end{aligned}
    $$
\end{bemerkungbox}

\begin{lemma}{}
    Sei $\{a(k)\}_{k=1}^n$ eine Orthonormalbasis von $\K^n$. Dann besitzt jeder Vektor $x\in \K^n$ die Darstellung
    \[
        x=\sum_{k=1}^n (x,a(k))\,a(k),
    \]
    und es gilt die Parseval-Gleichung
    \[
        \|x\|^2=\sum_{k=1}^n |(x,a(k))|^2.
    \]
\end{lemma}

\begin{satz}{Gram--Schmidt}
    Sei $\{a(1),\dots,a(n)\}$ eine Basis von $\K^n$. Dann entsteht durch
    \[
        b(1)=\frac{a(1)}{\|a(1)\|},\qquad
        \tilde{b}(k)=a(k)-\sum_{j=1}^{k-1}(a(k),b(j))\,b(j),\qquad
        b(k)=\frac{\tilde{b}(k)}{\|\tilde{b}(k)\|}
    \]
    eine Orthonormalbasis $\{b(1),\dots,b(n)\}$.
\end{satz}
